\documentclass[../DoAn.tex]{subfiles}
\begin{document}

\begin{center}
    \Large{\textbf{ABSTRACT}}\\
\end{center}
\vspace{1cm}
Air pollution is one of the most critical public health threats worldwide, and is particularly severe in Vietnam, where the annual average PM2.5 concentration is nearly six times higher than the WHO guideline. The national air quality monitoring network operates only approximately 35 automated stations, which is insufficient to capture micro-level air quality variations across urban areas. Existing commercial solutions such as PAM Air, IQAir, and PurpleAir all require WiFi or cellular connectivity at each monitoring point, incur high hardware costs, and use closed-source software. None of these systems simultaneously satisfy four criteria: multi-parameter sensing, long-range data transmission without WiFi, open-source availability, and low hardware cost. To address these limitations, this thesis adopts a four-tier IoT architecture combining low-cost sensors with LoRa long-range communication, eliminating the need for network infrastructure at each monitoring point.

The system consists of Sensor Nodes (ESP32 microcontroller integrated with PMS7003, CCS811, and AHT10 sensors) that measure six environmental parameters, encode the data into 18-byte binary packets, and transmit them via LoRa to a Gateway. The Gateway forwards the data over WiFi to a Backend Server (Node.js, TimescaleDB), which computes the Air Quality Index (AQI) according to US EPA standards and broadcasts real-time updates via Socket.IO to a React web frontend. The thesis presents four main contributions: (i) a fully operational end-to-end multi-tier IoT system comprising 13,759 lines of code deployed in a production environment, (ii) an 18-byte binary LoRa protocol that reduces air-time by approximately 11 times compared to JSON, coupled with a data deduplication mechanism, (iii) a telemetry processing pipeline on the Backend that ensures zero data loss, validated by 44 out of 44 test scenarios passing, and (iv) a hardware cost of approximately 895,000 VND (~36 USD) per Sensor Node --- roughly 11 times lower than the IQAir solution --- with the entire source code released as open-source.

\end{document}